\section{An existential condition excluding odd poles}\label{par:section}

We prove Proposition~\ref{par:pole-parity}, used in the finite rational
tests of Section~\ref{sec:reduction}.  It requires a positive existential
condition that holds on all integers and restricts the denominators of any
rational number satisfying it.  Square classes below mean classes in the
multiplicative group modulo squares.

Earlier existential formulas make related valuation information
accessible on restricted sets of primes.  For a fixed global field
\(K\) of characteristic different from \(2\) and a fixed quadratic
extension \(L/K\), Demeyer and Van Geel give an existential formula
which, for nonzero \(x,y\in K\), is equivalent to
\[
 v_{\mathfrak p}(x)\ \text{odd}
 \quad\Longrightarrow\quad
 v_{\mathfrak p}(x/y^2)>0
 \qquad
 \text{for every nonarchimedean prime \(\mathfrak p\) inert in \(L/K\)}.
\]
They also permit a fixed finite union of such inert-prime sets
\cite[Theorem~14 and Corollary~15 in the author version]{DemeyerVanGeel2006}.
Cornelissen and Zahidi use related odd-valuation predicates in their
elliptic-divisibility approach
\cite[Sections~3.3 and~3.11]{CornelissenZahidi2007}.
The existential definitions of Demeyer and Van Geel control the
chosen inert primes.
Proposition~\ref{par:pole-parity} instead supplies the stated
one-sided restriction outside a fixed finite exceptional set.

We first define a formula from finite multiplier lists and prove the
pole restriction. We then choose the lists once and for all and prove
integer completeness by realizing the same square class on two elliptic
curves.

\subsection{The formula and the pole restriction}

For now, let $\mathcal M\subset\Q^\times$ be a finite nonempty list,
and for each $m\in\mathcal M$ let $\mathcal D_m\subset\Q^\times$
be a finite nonempty list. The construction below associates a formula
$\Phi$ to these lists. Let $S_1$ be a finite set containing $2$, $3$,
and the supports of every member of all the lists. We will make one
fixed choice of the lists in the next subsection.

For $l\in\Q^\times$, write
\begin{equation}\label{par:curve}
 E_l:\quad y^2=x(x-l)(x+3l).
\end{equation}
The three roots are distinct, so this is an elliptic curve with its
point at infinity $O$ as origin. Define $\Phi(b)$ to hold if $b=0$,
or if $b\ne0$ and, for some $m\in\mathcal M$, $m_d\in\mathcal D_m$,
and $t,e\in\Q$, the following conditions hold:
\begin{equation}\label{par:auxiliary-functions}
 d=m_db,\qquad b'=mb,\qquad
 U=b\frac{1-dt^2}{1+dt^2},\qquad H'=b'(b'-U^2),
\end{equation}
\[
 1+dt^2\ne0,\qquad eH'\ne0,
\]
and each of $E_e$ and $E_{eH'}$ has a rational point whose
$x$-coordinate belongs to $b'\Q^{\times2}$.

For each fixed choice of the lists, this is a positive existential
formula. The choices of multipliers are finite disjunctions. A nonzero
condition $a\ne0$ is the existential equation $az=1$; the condition on
$x$ is expressed by $x=b'z^2$ with $z\ne0$. Introduce $U$ as a variable
and multiply its defining equality by $1+dt^2$. All remaining conditions
are polynomial equalities with fixed rational coefficients, which can be
cleared to integer coefficients.

\begin{lemma}\label{par:soundness}
For any finite lists as above, if $b\in\Q^\times$ satisfies $\Phi(b)$
and $q\notin S_1$, then $v_q(b)<0$ implies $v_q(b)\in2\Z$.
\end{lemma}

\begin{proof}
Suppose that $q\notin S_1$, that witnesses for $\Phi(b)$ have been
chosen, and that $k=v_q(b)<0$ is odd. The multipliers are $q$-adic
units, so $v_q(d)=v_q(b')=k$. If $t=0$, then $U=b$. If $t\ne0$,
the valuation $v_q(dt^2)=k+2v_q(t)$ is odd and hence nonzero.
When this valuation is positive, both $1-dt^2$ and $1+dt^2$ are
units. When it is negative, both have valuation $v_q(dt^2)$.
In either case $v_q(U)=k$. Since $2k<k$,
\[
 v_q(b'-U^2)=2k,\qquad v_q(H')=3k.
\]
Consequently $v_q(e)$ and $v_q(eH')$ have opposite parities. Let
$l$ be the one of these two parameters whose valuation is even,
say $v_q(l)=2a$. Replacing
\[
 l=q^{2a}L,\qquad x=q^{2a}X,\qquad y=q^{3a}Y
\]
gives $Y^2=X(X-L)(X+3L)$ with $L\in\Z_q^\times$ and does not
change the square class of $x$.

Every nonzero $X$ on this unit-parameter curve has even valuation.
Indeed, if $v_q(X)<0$, all three factors have valuation $v_q(X)$,
so the right side has valuation $3v_q(X)$. If $v_q(X)>0$, the
other two factors are units because $q\ne3$, so its valuation is
$v_q(X)$. In both cases this valuation must be even. The case
$v_q(X)=0$ already has the desired parity. This contradicts
$x\in b'\Q_q^{\times2}$, since $v_q(b')=k$ is odd.
\end{proof}


\subsection{Preparing integer witnesses}

We now choose the multiplier lists so that the formula also holds on
every integer. These lists will be fixed independently of the input
$b$. Choose a finite list $\mathcal M\subset\Q^\times$ with the
following property: for every $b\in\Q^\times$, some $m\in\mathcal M$
makes
\begin{equation}\label{par:b-prime-choices}
 b'=mb>0,\qquad b'\in\Q_2^{\times2},\qquad
 b'\in\text{the unit nonsquare class of }\Q_3^\times.
\end{equation}
Here the last condition means that $b'$ has even $3$-adic valuation
and nonsquare unit part. Such a list exists because
$\R^\times/\R^{\times2}$, $\Q_2^\times/\Q_2^{\times2}$, and
$\Q_3^\times/\Q_3^{\times2}$ are finite, their classes are open,
and weak approximation supplies a rational representative for each
specified finite tuple of classes. For each $m\in\mathcal M$, set
\[
 S(m)=\{2,3\}\cup\operatorname{Supp}(m).
\]
The same argument gives a finite list
$\mathcal D_m\subset\Q^\times$ such that, for any $b\ne0$, some
$m_d\in\mathcal D_m$ makes
\begin{equation}\label{par:d-choices}
 d=m_db>0,\qquad d\in\Q_q^{\times2}\quad(q\in S(m)).
\end{equation}
Use these fixed lists in $\Phi$, and fix $S_1$ containing $2$, $3$,
and all their supports as above. The normalizations
\eqref{par:b-prime-choices}--\eqref{par:d-choices} are used only to
select entries $m,m_d$ for an integer input; they are not additional
conjuncts of $\Phi$. Lemma~\ref{par:soundness} already proves the pole
restriction for this choice.

Fix $b\in\Z\setminus\{0\}$ and choose $m,m_d$ so that
\eqref{par:b-prime-choices} and \eqref{par:d-choices} hold.
We choose $t$ to serve two purposes. The value $U$ should be small at
the places where $d$ was made a local square, while at primes outside
$S(m)$ where $v_q(b')$ is odd it should retain the valuation of $b$.
These properties will make $H'$ a local square at every prime where
$v_q(b')$ is odd, and $b'$ a local square at every prime where
$v_q(H')$ is odd.

At the real place and at every $q\in S(m)$, choose a local solution
of $dt^2=1$. As $t$ tends to this solution, the rational function
$U$ tends to zero: its denominator tends to the nonzero number $2$.
Thus, in sufficiently small neighborhoods, $H'>0$ at the real place,
and
\[
 H'=b'^2\left(1-\frac{U^2}{b'}\right)\in\Q_q^{\times2}
 \quad(q\in S(m)),
\]
by openness of the local square subgroup. At every
$q\in\operatorname{Supp}(m_d)\setminus S(m)$, require $t$ to be
so small that $v_q(dt^2)>0$. Weak approximation satisfies all these
conditions simultaneously with a rational $t$. In particular
$1+dt^2\ne0$ and $H'>0$.

Let $\beta,\theta$ be the unique positive squarefree integers
representing the square classes of $b',H'$, and let
\[
 B=\operatorname{Supp}(\beta),\qquad
 J=\operatorname{Supp}(\theta).
\]

\begin{lemma}\label{par:cross-squares}
The sets $B,J$ are disjoint and avoid $2,3$. Moreover,
\begin{equation}\label{par:mutual-squares}
 \theta\in\Q_q^{\times2}\quad(q\in B),\qquad
 \beta\in\Q_q^{\times2}\quad(q\in J).
\end{equation}
Both $\beta$ and $\theta$ are squares in $\Q_2^\times$; at $3$,
$\theta$ is a square and $\beta$ is a unit nonsquare. In particular
$B\ne\varnothing$.
\end{lemma}

\begin{proof}
The assertions at $2,3$ follow from the choices already made and
are unchanged on replacing numbers by square-class representatives.
At a prime $q\in B\cap S(m)$, $H'$ is a local square by construction.
Now suppose $q\in B\setminus S(m)$. Since $m$ is a unit at $q$ and
$b$ is an integer, $v_q(b)=v_q(b')$ is a positive odd integer.
If $q\notin\operatorname{Supp}(m_d)$ and $t\ne0$, then
$v_q(dt^2)=v_q(b)+2v_q(t)$ is odd and nonzero, and the quotient
$(1-dt^2)/(1+dt^2)$ is a unit, just as in
Lemma~\ref{par:soundness}. For $t=0$ this quotient is $1$.
If instead $q\in\operatorname{Supp}(m_d)$, our additional smallness
condition on $t$ also makes it a unit. Hence $v_q(U)=v_q(b)$, and
\[
 v_q(U^2/b')=v_q(b)>0.
\]
As $q$ is odd, $1-U^2/b'$ is a local square. The displayed formula
for $H'$ proves the first assertion of \eqref{par:mutual-squares},
and also proves $B\cap J=\varnothing$.

For the converse assertion, let $q\in J$. Such a prime is outside
$S(m)$, since $H'$ is square there, and $v_q(b')$ is even because
$q\notin B$. If $U=0$, then $H'=b'^2$ and there is no such prime,
so suppose $U\ne0$. Put
\[
 \alpha=v_q(b'),\qquad \gamma=2v_q(U).
\]
Both are even. If $\alpha\ne\gamma$, then
$v_q(b'-U^2)=\min(\alpha,\gamma)$ is even, contrary to
the oddness of $v_q(H')$. Therefore $\alpha=\gamma$.
Since $v_q(H')=\alpha+v_q(b'-U^2)$ is odd, cancellation is
strict: $v_q(b'-U^2)>\alpha$. Thus
$b'/U^2\in1+q\Z_q$, a square in $\Q_q^\times$, proving that
$\beta$ is square there. This argument also covers possible
denominators in $t$ or $U$. Finally $\beta$ cannot be $1$, since
its class at $3$ is nonsquare.
\end{proof}


Write $H'=\theta c^2$ with $c\in\Q^\times$. For every $n\in\Q^\times$
the isomorphism
\begin{equation}\label{par:twist-scaling}
 E_{n\theta}\longrightarrow E_{nH'},\qquad
 (x,y)\longmapsto(c^2x,c^3y)
\end{equation}
preserves $x$-coordinate square classes. Thus, taking $e=n$ in the
formula, it remains to find one $n$ for which both $E_n$ and
$E_{n\theta}$ have a rational point with nonzero $x$-coordinate of class
$\beta$, the class of $b'$. We first give a criterion for one parameter
$l$, and then construct a single $n$ for which it applies to both twists.

\subsection{A criterion for one twist}

We recall the exact part of full two-descent that we use; see
\cite[Chapter~X, Proposition~1.4, p.~315]{Silverman2009}. For a split
separable cubic $y^2=\prod_{i=1}^3(x-e_i)$ over a field $k$ of
characteristic zero, choosing two roots identifies
$H^1(k,E[2])$ with $(k^\times/k^{\times2})^2$. In the coordinate
corresponding to $e_i$, the Kummer map takes a point to the class
of $x-e_i$; at $(e_i,0)$ the value is replaced by
$\prod_{j\ne i}(e_i-e_j)$, and at $O$ it is $1$.
The Kummer map is an injection of $E(k)/2E(k)$ into this group.
The group $\Sel_2(E/\Q)$ is the subgroup of rational pairs lying
in the Kummer image at every completion. The standard Selmer exact
sequence \cite[Chapter~X, Theorem~4.2, p.~333]{Silverman2009} is
\begin{equation}\label{par:selmer-exact-sequence}
 0\longrightarrow E(\Q)/2E(\Q)
 \longrightarrow\Sel_2(E/\Q)
 \longrightarrow\Sha(E/\Q)[2]\longrightarrow0.
\end{equation}


We use the coordinates $(x+3l,x-l)$ for $E_l$. For a finite rational
point $T$ with $x(T)\notin\{0,l,-3l\}$, the curve equation gives
\[
             [x(T)]=[x(T)+3l]\,[x(T)-l].
\]
Consequently, if the class $(\beta,1)$ is represented by a rational
point and is distinct from the classes of all four rational two-torsion
points, that point has the required nonzero $x$-coordinate class $\beta$.
The next calculation identifies conditions that make $(\beta,1)$ the
one additional Selmer direction beyond rational two-torsion.

For the Selmer graph-matrix viewpoint underlying this calculation, see
Monsky's appendix to Heath-Brown \cite[Appendix, pp.~365--370]{HeathBrown1994}.
The matrix and simultaneous rank calculations below adapt that viewpoint
to the present curve family.

At the primes dividing $l$, the local square conditions will be recorded
by binary residue symbols. For an odd prime $q$ and a $q$-adic unit $a$,
write $[a/q]\in\F_2$ for the nonsquare bit: it is $0$ for a square
residue and $1$ for a nonsquare residue. For a finite set $Q$ of distinct
primes different from $2,3$, indexed by $i$, and a vector $w\in\F_2^Q$,
write $D_w$ for the diagonal matrix with diagonal $w$. Put
\begin{equation}\label{par:matrices}
 s_i=[-1/q_i],\qquad \lambda_i=[3/q_i],\qquad
 (Ku)_i=\sum_{j\ne i}[q_j/q_i](u_j+u_i),\qquad N=K+D_s.
\end{equation}
All matrices and vectors in this calculation are over $\F_2$. Write
$\mathbf1$ for the all-ones vector and $\mathbf1_B$ for the indicator
of $B$ in $Q$. Quadratic reciprocity gives
\begin{equation}\label{par:reciprocity-matrix}
 K+K^{\mathsf t}=ss^{\mathsf t}+D_s.
\end{equation}
In particular, if $s^{\mathsf t}\mathbf1=1$, then
\begin{equation}\label{par:matrix-null-vectors}
 K\mathbf1=K^{\mathsf t}\mathbf1=0,
 \qquad N=K^{\mathsf t}+ss^{\mathsf t}.
\end{equation}

The family \(y^2=x(x-n)(x+3n)\) also occurs in the study of
\(\pi/3\)- and \(2\pi/3\)-congruent numbers.  Mokrani adapted
Monsky matrices to these families \cite{Mokrani2020}.  For this
family, Wei and Guo give a two-Selmer matrix formulation and write
the same quadratic-reciprocity identity for the associated R\'edei
matrix as \eqref{par:reciprocity-matrix}
\cite[Section~4 and Theorem~6.2]{WeiGuo2022}.  The simultaneous
kernel conditions for the two twists required here are constructed
below.

\begin{lemma}\label{par:descent}
Let $l>0$ be squarefree, prime to $6$, with
$l\equiv3\pmod4$ and $l\equiv2\pmod3$. Suppose that its support
$Q$ contains the nonempty set $B$ of Lemma~\ref{par:cross-squares}
and a prime outside $B$, and suppose that
\begin{equation}\label{par:matrix-target}
 \ker K=\langle\mathbf1,\mathbf1_B\rangle,
 \qquad \lambda\notin\operatorname{im}K.
\end{equation}
Then
\[
 \dim_{\F_2}\Sel_2(E_l/\Q)=3,
\]
and $(\beta,1)$ is a Selmer class outside the subgroup of rational
two-torsion classes, in Kummer coordinates $(x+3l,x-l)$.
\end{lemma}

\begin{proof}
In the stated coordinates, the points $(l,0)$ and $(-3l,0)$ have
classes
\begin{equation}\label{par:torsion-classes}
 \tau_1=(l,1),\qquad \tau_2=(3,-l),
\end{equation}
respectively. For example, the replacement coordinate at $(l,0)$
is $(l+3l)l=4l^2$, and the replacement at $(-3l,0)$ is
$(-3l)(-4l)=12l^2$. These two classes are independent globally:
$l$ is nonsquare, and the second coordinate of $\tau_2$ is negative.

We first bound the Selmer group using necessary local conditions.
Afterward we verify that $(\beta,1)$ lies in the actual local Kummer
image everywhere. For any odd prime $q$, the group
$E_l(\Q_q)/2E_l(\Q_q)$ has order four. To see this, choose an
open formal subgroup $H$ on which multiplication by $2$ is an
isomorphism. The quotient $G=E_l(\Q_q)/H$ is finite, and
$E_l(\Q_q)/2E_l(\Q_q)\simeq G/2G$. Moreover
$E_l(\Q_q)[2]\simeq G[2]$: injectivity follows from $H[2]=0$,
and a lift of any element of $G[2]$ can be corrected by an element
of $H$, using the surjectivity of multiplication by $2$ on $H$.
Finally $|G/2G|=|G[2]|=4$.

At an odd prime $q\nmid3l$, the roots are integral and pairwise
distinct modulo $q$. If $v_q(x)<0$, all three factors have that
valuation, which must be even. If $x$ is integral, at most one
factor has positive valuation, and that valuation must be even.
The replacement classes at the roots are units as well. Hence
the local Kummer image is contained in the subgroup of pairs
having even coordinate valuations. This unramified subgroup has
order four, so the containment is equality.

At a support prime $q_i\mid l$, the valuation pairs of $\tau_1$
and $\tau_2$ are $(1,0)$ and $(0,1)$ modulo $2$; therefore they
generate the entire local image, again by its order four. At $3$,
$l$ is a unit nonsquare and $-l$ is a square. Thus $\tau_1$
and $\tau_2$ have independent first-coordinate classes and square
second coordinates. The local image at $3$ is consequently
\begin{equation}\label{par:local-three}
 (\Q_3^\times/\Q_3^{\times2})\times\{1\}.
\end{equation}

At $2$ we need the necessary condition that both Kummer coordinates
have even valuation. For a point away from the roots with
$v_2(x)<0$, all three factors have valuation $v_2(x)$, so that
valuation is even. If $x$ is integral and even, $x+3l$ and $x-l$
are units. Suppose $x$ is odd. If $x-l=2u$ with $u$ odd, then
\[
 (y/2)^2=(l+2u)u(u+2l)\equiv l\pmod8.
\]
The right side is an odd unit congruent to $3$ or $7$ modulo $8$,
which is impossible for a square. In every remaining nonroot
case, $v_2(x-l)\ge2$ and $v_2(x+3l)\ge2$. Their difference is
$4l$, so their minimum valuation is $2$. Since $x$ is a unit,
their valuation sum is even, and hence both valuations are even.
The root and identity classes also have even valuations, as follows
from \eqref{par:torsion-classes}. This establishes the claimed
necessary condition at $2$; no sufficiency is being asserted.

At the real place the first coordinate is always positive. Indeed,
for a nonroot real point the allowable $x$ intervals are
$(-3l,0)$ and $(l,\infty)$, and the replacement first coordinates
at the roots are positive as well. The sign of the second coordinate
therefore gives a homomorphism from the Selmer group to
$\{\pm1\}$, and it is surjective because of $\tau_2$. Its kernel,
the classes with both coordinates positive, has index two.
Every class in this kernel has a unique representative of the form
\begin{equation}\label{par:normalized-pair}
 \left(3^\delta\prod_iq_i^{u_i},\ \prod_iq_i^{v_i}\right),
 \qquad \delta\in\F_2,\quad u,v\in\F_2^Q.
\end{equation}
Here the good-prime restrictions and the even valuations at $2$
exclude every other prime, and \eqref{par:local-three} excludes
$3$ from the second coordinate.

At $q_i$, compare \eqref{par:normalized-pair} with the local
torsion class
\[
 \tau_1^{u_i}\tau_2^{v_i}
   =(l^{u_i}3^{v_i},(-l)^{v_i}).
\]
The coordinate valuations already agree modulo $2$. Requiring
the resulting unit quotients to be squares gives exactly
\begin{equation}\label{par:selmer-matrix-equations}
 Ku+D_\lambda v=\delta\lambda,\qquad Nv=0.
\end{equation}
For example the first unit quotient has nonsquare bit
$\sum_{j\ne i}[q_j/q_i](u_j+u_i)+(\delta+v_i)\lambda_i$;
the second has bit $(Kv)_i+s_iv_i$. The requirement that the
second coordinate be square at $3$ adds
\begin{equation}\label{par:selmer-three-equation}
 (\lambda+s)^{\mathsf t}v=0,
\end{equation}
using $[q_i/3]=\lambda_i+s_i$.

The residue assumptions imply
$s^{\mathsf t}\mathbf1=1$ and $\lambda^{\mathsf t}\mathbf1=0$.
Equations \eqref{par:matrix-null-vectors} give
\begin{equation}\label{par:kernel-N}
 \ker N=\ker K^{\mathsf t}\cap s^\perp.
\end{equation}
Indeed, left multiplying $Nv=0$ by $\mathbf1^{\mathsf t}$
gives $s^{\mathsf t}v=0$, and the identity
$N=K^{\mathsf t}+ss^{\mathsf t}$ then gives the claim in both
directions. By \eqref{par:matrix-target}, $\ker K^{\mathsf t}$
has dimension two. It contains $\mathbf1$, on which
$s^{\mathsf t}$ is nonzero, so \eqref{par:kernel-N} is a line.
The functional $\lambda^{\mathsf t}$ is nonzero on
$\ker K^{\mathsf t}$ because
$\lambda\notin\operatorname{im}K=(\ker K^{\mathsf t})^\perp$.
Since it vanishes on $\mathbf1$, it is nonzero on the complementary
line \eqref{par:kernel-N}. Equations
\eqref{par:selmer-three-equation} and \eqref{par:kernel-N}
therefore force $v=0$. The first equation in
\eqref{par:selmer-matrix-equations} now gives
$Ku=\delta\lambda$, so $\delta=0$ and
$u\in\ker K=\langle\mathbf1,\mathbf1_B\rangle$.
There are at most four sign-normalized classes, and hence at
most eight Selmer classes.

For equality, we verify that
\[
 \gamma=(\beta,1)
\]
is an actual local Kummer class everywhere. At $2$ and at the real
place it is the identity class, since $\beta$ is a positive local
square. At $3$ it is $\tau_1$, since $\beta$ and $l$ are both
unit nonsquares. At $q_i\mid l$, take $u=\mathbf1_B$ and $v=0$:
the comparison with $\tau_1^{u_i}$ has square coordinate quotients
precisely because $K\mathbf1_B=0$. At every other finite odd
prime it is unramified, so lies in the local image already
described. This proves $\gamma\in\Sel_2(E_l/\Q)$.
It is neither the identity nor $\tau_1$: $B$ is nonempty, and
$l$ has a prime factor outside $B$. It cannot equal $\tau_2$ or
$\tau_1\tau_2$, whose second coordinates are negative. Thus
$\tau_1,\tau_2,\gamma$ are three independent Selmer classes,
and the upper bound is attained.
\end{proof}


It remains to realize this particular Selmer class by a rational point.
The consequence we use from Theorem~1.1 of \cite{PointwiseTwoConverse}
is the following: an elliptic curve $C/\Q$ with $C(\Q)[2]\ne0$ and
$\corank_{\Z_2}\Sel_{2^\infty}(C/\Q)\in\{0,1\}$ has finite
$\Sha(C/\Q)$. The next proof verifies the full Selmer-corank hypothesis
before applying it.

\begin{lemma}\label{par:global-point}
Under the hypotheses of Lemma~\ref{par:descent}, the curve $E_l$
has a rational point with nonzero $x$-coordinate in
$\beta\Q^{\times2}$.
\end{lemma}

\begin{proof}
Put $r=\rank E_l(\Q)$. Since $E_l$ has full rational two-torsion,
the Kummer sequence \eqref{par:selmer-exact-sequence} and
Lemma~\ref{par:descent} imply
\begin{equation}\label{par:rank-sha}
 r+\dim_{\F_2}\Sha(E_l/\Q)[2]=1.
\end{equation}
This also controls the full $2$-primary Selmer group before any
finiteness conclusion. The usual exact sequence is
\[
 0\longrightarrow E_l(\Q)\otimes\Q_2/\Z_2
 \longrightarrow\Sel_{2^\infty}(E_l/\Q)
 \longrightarrow\Sha(E_l/\Q)[2^\infty]\longrightarrow0.
\]
These are cofinitely generated $\Z_2$-modules. If the last group
has corank $t$, its structure is
$(\Q_2/\Z_2)^t\oplus F$ with $F$ finite, so
$t\le\dim_{\F_2}\Sha(E_l/\Q)[2]$. It follows from
\eqref{par:rank-sha} that
\[
 0\le\corank_{\Z_2}\Sel_{2^\infty}(E_l/\Q)=r+t\le1.
\]
The curve has nonzero rational two-torsion, so
\cite[Theorem~1.1]{PointwiseTwoConverse} now applies and proves
that $\Sha(E_l/\Q)$ is finite.

The Cassels--Tate pairing on this finite group is perfect and
alternating. Alternation here uses the principal polarization
represented by the rational divisor $(O)$; see
\cite{Cassels1962}; see also
\cite[Section~1 and Corollary~9]{PoonenStoll1999}. For completeness,
an alternating perfect pairing on a finite abelian $2$-group
forces an even number of cyclic factors. Choose an element $x$
of maximal order $2^a$. Perfection supplies $y$ whose pairing
with $x$ has exact order $2^a$. The subgroup generated by $x,y$
is a nondegenerate copy of $(\Z/2^a\Z)^2$, and the whole group
is its direct sum with its orthogonal complement. Induction gives
a decomposition into such paired cyclic factors. Hence
\[
 \dim_{\F_2}\Sha(E_l/\Q)[2]\quad\text{is even}.
\]
This conclusion uses the perfect pairing on the full finite
$2$-primary group; its restriction to the subgroup killed by $2$
need not be perfect. Equation~\eqref{par:rank-sha} therefore gives
$\Sha(E_l/\Q)[2]=0$ and $r=1$.

The Kummer sequence now shows that $\gamma=(\beta,1)$ is the
class of a rational point $T$. By Lemma~\ref{par:descent} its class
differs from all four two-torsion point classes. Thus $T$ is finite
and none of $x(T)$, $x(T)-l$, $x(T)+3l$ is zero. The curve equation
gives, modulo squares,
\[
 [x(T)]=[x(T)+3l]\,[x(T)-l]=\beta,
\]
as required.
\end{proof}


\begin{remark}
The unrestricted $2$-converse \cite[Theorem~1.1]{UnrestrictedTwoConverse}
can replace the pointwise theorem in the preceding proof. It gives
finiteness of the whole $\Sha(E_l/\Q)$ from the same full Selmer-corank
bound, without any rational-two-torsion hypothesis.
\end{remark}

\subsection{One parameter for both twists}

The criterion is now a condition on the prime support of one twist.
We construct a single positive squarefree $n$ for which it holds both
on the support of $n$ and on that of $n\theta$. The mutual local-square
relations from Lemma~\ref{par:cross-squares} are what allow the second
support to be added without disturbing the first construction.

\begin{lemma}\label{par:simultaneous-matrices}
For $\beta,\theta,B,J$ as in Lemma~\ref{par:cross-squares}, there is
a positive squarefree integer $n$, prime to $6\theta$, whose support
$P$ contains $B$ and at least one auxiliary prime outside $B$, such
that the following statements hold for both $l=n$ and $l=n\theta$.
The number $l$ satisfies $l\equiv3\pmod4$ and $l\equiv2\pmod3$.
On its prime support $Q=P$ or $Q=P\cup J$, the matrices
\eqref{par:matrices} satisfy \eqref{par:matrix-target}.
\end{lemma}

\begin{proof}
Since $\beta,\theta$ are odd squares at $2$, their products of
$(-1)$-symbols are $1$. At $3$, reciprocity gives
$[q/3]=[3/q]+[-1/q]=\lambda_q+s_q$. Consequently
\begin{equation}\label{par:fixed-symbol-sums}
 \sum_{i\in B}s_i=\sum_{j\in J}s_j=0,\qquad
 \sum_{i\in B}\lambda_i=1,\qquad
 \sum_{j\in J}\lambda_j=0.
\end{equation}
The fixed $B$--$J$ symbols need not vanish individually. Their
row sums, however, satisfy
\begin{equation}\label{par:cross-symbol-sums}
 \sum_{j\in J}[q_j/q_i]=0\quad(i\in B),\qquad
 \sum_{i\in B}[q_i/q_j]=0\quad(j\in J),
\end{equation}
by \eqref{par:mutual-squares}.

Introduce one auxiliary prime $q_0$ of type
$(s_0,\lambda_0)=(1,1)$, and a set $I$ of $|B|-1$ auxiliary
primes of type $(0,0)$. These types can be imposed by the residue
classes $7$ and $1$ modulo $12$, respectively. We first prescribe
all required Legendre symbols abstractly, and realize them by actual
primes at the end. Until then, each unchosen symbol is a bit whose reverse
is constrained by quadratic reciprocity; the matrix reciprocity identity
therefore holds for every completion of these choices. On each of the sets
\[
 Q_0=B\cup I\cup\{q_0\},\qquad Q_1=Q_0\cup J,
\]
the total $s$-sum is $1$ and the total $\lambda$-sum is $0$.
The former gives $l\equiv3\pmod4$, while the sum of $s+\lambda$
gives $l\equiv2\pmod3$ for the corresponding prime product.

We first show that, subject to $K\mathbf1_B=0$, both matrix targets follow
from one nonsingularity condition. Fix $b_*\in B$ and let $a$ have entries
$1$ at $q_0,b_*$ and $0$ elsewhere. Thus
\[
 a^{\mathsf t}\mathbf1=0,\qquad
 a^{\mathsf t}\mathbf1_B=1.
\]
We shall impose $K\mathbf1_B=0$ on both supports. Assuming this condition
for the moment, set $K^\#=K+\lambda a^{\mathsf t}$. The total symbol
sums and \eqref{par:matrix-null-vectors} show that $K^\#$ kills
$\mathbf1$ on both sides, and
\begin{equation}\label{par:sharp-B}
 K^\#\mathbf1_B=\lambda.
\end{equation}
Put
\[
 T=\operatorname{Id}-\mathbf1_Ba^{\mathsf t}.
\]
This is a projection onto $a^\perp$, with kernel
$\langle\mathbf1_B\rangle$, and it fixes $\mathbf1$. By
\eqref{par:sharp-B},
\begin{equation}\label{par:projection-factorization}
 K=K^\#T.
\end{equation}
Suppose that $\ker K^\#=\langle\mathbf1\rangle$. If $Kx=0$, then
$Tx=c\mathbf1$, so
$x=c\mathbf1+(a^{\mathsf t}x)\mathbf1_B$; the converse follows from
$K\mathbf1=K\mathbf1_B=0$. This gives the required kernel of $K$.
If $Kx=\lambda$, put $y=Tx$. Then $a^{\mathsf t}y=0$ and
$K^\#(y-\mathbf1_B)=0$, so
$y-\mathbf1_B\in\langle\mathbf1\rangle$. Applying $a^{\mathsf t}$
would give $-1=0$, a contradiction. Thus $\lambda\notin\operatorname{im}K$.

For either support, let $M$ be the matrix obtained from $K^\#$ by
deleting the $q_0$ row and column. It is enough to make $M$ nonsingular.
Indeed, if $K^\#x=0$, subtract $x_{q_0}\mathbf1$ from $x$. The resulting
vector still lies in the kernel, has zero $q_0$ coordinate, and its
remaining coordinates are killed by $M$. Nonsingularity makes that
vector zero, proving $\ker K^\#=\langle\mathbf1\rangle$. We therefore
have two tasks for the symbol choices: impose $K\mathbf1_B=0$ and make
these deleted matrices nonsingular on both supports.

For the smaller support, let $C$ be a $B$-by-$I$ matrix whose columns
are a basis of the even-sum hyperplane
\[
 \{x\in\F_2^B:\mathbf1_B^{\mathsf t}x=0\}.
\]
These columns will form the off-diagonal block of the deleted matrix.
When $|B|=1$, $C$ is the empty matrix. Prescribe
$[q_i/q_b]=[q_b/q_i]=C_{bi}$ for $b\in B,i\in I$.
Every other off-diagonal symbol incident with $I$ is set to zero.
Prescribe
\[
 [q_0/q_b]=\sum_{i\in I}C_{bi}\qquad(b\in B),
\]
and obtain reverse symbols by reciprocity. In particular
$[q_b/q_0]=[q_0/q_b]+s_b$. The $q_0$--$J$ symbols remain free.

These prescriptions ensure $K\mathbf1_B=0$ on both supports.
At a $B$ row, the terms from $I$ and $q_0$ cancel, and the sum
from $J$, when present, is zero by \eqref{par:cross-symbol-sums}.
At an $I$ row the sum is a column sum of $C$, hence zero. At the
$q_0$ row it is the sum of all entries of $C$ plus $\sum_B s_b$,
again zero. At a $J$ row it is zero by the second equality in
\eqref{par:cross-symbol-sums}. It therefore remains to make the deleted
matrices nonsingular.

On $Q_0$, the deleted matrix has the form
\begin{equation}\label{par:block}
 M_0=\begin{pmatrix}A&C\\ C^{\mathsf t}&0\end{pmatrix},
 \qquad \mathbf1_B^{\mathsf t}A\mathbf1_B=1.
\end{equation}
The equality follows from \eqref{par:sharp-B} and
$\sum_B\lambda_b=1$. The lower right block is zero since the
only potentially nonzero symbols incident with $I$ are in $C$,
whose column sums vanish. The added matrix $\lambda a^{\mathsf t}$
changes neither off-diagonal block, because $a$ and $\lambda$
both vanish on $I$. No symmetry of $A$ is required.

To prove $M_0$ nonsingular, let $M_0(x,y)^{\mathsf t}=0$.
The second block equation says $C^{\mathsf t}x=0$. Since the
columns of $C$ span the even-sum hyperplane, its annihilator is
$\langle\mathbf1_B\rangle$; hence $x=c\mathbf1_B$.
Left multiplying the first block equation
$cA\mathbf1_B+Cy=0$ by $\mathbf1_B^{\mathsf t}$ gives $c=0$
by \eqref{par:block}. The injectivity of $C$ then gives $y=0$.
This proof applies also when $|B|$ is even; when $|B|=1$ it says
simply that $M_0=(1)$.

On adding $J$, the $B,I$ block of the deleted matrix stays equal to
$M_0$. Indeed, the only possible changes to its diagonal are the
$B$--$J$ row sums, which vanish by \eqref{par:cross-symbol-sums}; all
$I$--$J$ symbols were set to zero. Now varying the bit $[q_0/q_j]$ for
one $j\in J$ also varies its reverse by the same bit, as their
reciprocity discrepancy is fixed. After deleting the $q_0$ row and
column, the only surviving change is a toggle of the $j$-th diagonal
entry. Thus these choices independently toggle the $J$-diagonal entries
of the enlarged matrix. Its determinant is a multilinear polynomial
in these bits whose coefficient of their full product is $\det M_0=1$.
A nonzero multilinear polynomial over $\F_2$ cannot vanish at every
point of the Boolean cube: this follows by induction on the number of
variables, writing it as $f_0+Xf_1$. Hence some choice makes the enlarged
matrix nonsingular. If $J$ is empty, the smaller-support calculation
already suffices. The chosen bits do not affect the smaller matrix.

The deleted matrices are now nonsingular on both supports, so the
preceding reduction proves \eqref{par:matrix-target} on both.
Finally realize the prescribed symbols. Choose the auxiliary primes
successively. At each step, specify the type modulo $12$ and the
Legendre symbols modulo every already fixed prime in $B\cup J$
or among the earlier auxiliary primes. Each symbol can be imposed
by a nonzero residue modulo that prime; reverse symbols are
consistent by the prescribed reciprocity rule. The Chinese
remainder theorem gives a reduced residue class, and Dirichlet's
theorem supplies infinitely many primes in it. Excluding the
finitely many previously used primes causes no difficulty.
The resulting product
$n=\prod_{q\in B\cup I\cup\{q_0\}}q$ has all the asserted
properties.
\end{proof}

\begin{proof}[Completion of the proof of Proposition~\ref{par:pole-parity}]
The value $b=0$ satisfies $\Phi$ by definition. For a nonzero integer
$b$, make the local choices preceding Lemma~\ref{par:cross-squares} and choose
$n$ by Lemma~\ref{par:simultaneous-matrices}. Both $l=n$ and
$l=n\theta$ satisfy the hypotheses of Lemma~\ref{par:global-point}, so
each curve has a point with nonzero $x$-coordinate of class $\beta$,
which is also the class of $b'$. The isomorphism
\eqref{par:twist-scaling} sends the point on $E_{n\theta}$ to one on
$E_{nH'}$ with the same $x$-coordinate square class. Taking $e=n$
gives the two points required in the definition of $\Phi$, on $E_e$
and $E_{eH'}$. This proves integer completeness, while
Lemma~\ref{par:soundness} proves the asserted pole restriction.
\end{proof}
